\section{函数与场的微分算子}
\label{sec:mcalc-operators}

欧氏空间中，梯度、Hessian、散度与Laplace算子都由偏导数组合而成，似乎只需要坐标。实际上它们各自悄悄用了欧氏结构：梯度要把$df$变成向量，用到了内积；Hessian要比较不同点处的梯度，用到了平行的坐标基；散度衡量体积的变化，用到了体积元。在流形上，这些结构分别由度量、Levi-Civita联络和黎曼体积提供。本节逐一定义这四个算子，并写明每个算子的输入、输出和所依赖的结构；它们之间的积分关系与边界项留到下一节。

\subsection{梯度}
\label{sec:mcalc-gradient}

函数$f$的微分$df_p$是余切向量，它告诉我们沿每个方向$v$的变化率$df_p(v)$，本身却不是一个方向。要回答“沿哪个方向增长最快”，需要把线性读数$df_p$表示成与某个向量的内积，而这只有在给定内积后才可能。

\begin{definition}[梯度]
\label{def:mcalc-gradient}
设$(M,g)$是黎曼流形，$f\in C^\infty(M)$。$f$的\term{梯度}（gradient）是唯一满足
\[
g\bigl(\operatorname{grad}f,\,X\bigr)=df(X)=Xf\qquad\text{对所有}X\in\mathfrak X(M)
\]
的向量场$\operatorname{grad}f$。
\end{definition}

每点处$g_p$是正定内积，线性泛函$df_p$由唯一的向量通过内积表示，所以梯度存在唯一。在坐标中，$df=\sum_j\partial_jf\,dx^j$，令$\operatorname{grad}f=\sum_ia^i\partial_i$，条件化为$\sum_ig_{i,j}a^i=\partial_jf$，于是
\begin{equation}
\operatorname{grad}f=\sum_{i,j}g^{i,j}\,\partial_jf\;\partial_i .
\label{eq:mcalc-gradient-coordinates}
\end{equation}
输入是函数，输出是向量场，用到度量而不用联络或体积。

\paragraph{最陡增加方向}
对单位切向量$u\in T_pM$，方向导数为
\begin{equation*}
df_p(u)=g_p(\operatorname{grad}f,u)\leq\lVert\operatorname{grad}f\rVert_g,
\end{equation*}
等号当且仅当$u$与$\operatorname{grad}f$同向，这是Cauchy--Schwarz不等式。所以$\operatorname{grad}f(p)\neq0$时，它的方向是$f$增加最快的单位方向，长度是最大增长率。“单位”是用$g$量出的，换了度量，单位向量的集合就变了，最陡方向也随之改变。

\paragraph{梯度依赖度量}
同一个函数在不同度量下有不同的梯度。在$\mathbb R^2$上取$f(x,y)=x+y$，欧氏度量下$\operatorname{grad}f=(1,1)$；若取常度量$\bm G=\operatorname{diag}(1,4)$，由式\eqref{eq:mcalc-gradient-coordinates}，$\operatorname{grad}f=(1,1/4)$。后一度量把$y$方向的单位长度放大一倍，同样的函数增量在$y$方向“更远”，最陡方向便偏向$x$轴。

即使度量不变、只换坐标，也不能把偏导数排成的列向量当作梯度。欧氏平面的极坐标中，
\[
\operatorname{grad}f=\partial_rf\,\partial_r+\frac1{r^2}\partial_\theta f\,\partial_\theta ,
\]
而$(\partial_rf,\partial_\theta f)$不是这个向量的分量。偏导数组成的是$df$的分量，它与梯度的分量相差一个$\bm G^{-1}$。对取诱导度量的嵌入子流形$M\subseteq\mathbb R^D$，若$f$是某个环境函数$\tilde f$的限制，则$\operatorname{grad}f=\bm P\,\nabla\tilde f$，即环境梯度在切空间上的正交投影：对切向量$X$，$\langle\bm P\nabla\tilde f,X\rangle=\langle\nabla\tilde f,X\rangle=Xf$。流形约束下的梯度方法正是从这里出发，可参考\BookChapterRef{chap:manifold-optimization}{卷五“流形约束下的优化”章}。

\subsection{Hessian}
\label{sec:mcalc-hessian}

欧氏空间中，Hessian矩阵$[\partial_i\partial_jf]$描述函数的二阶变化。在流形的坐标中直接取二阶偏导数会得到依赖坐标的结果：极坐标下$f=r^2$的二阶偏导数矩阵是$\operatorname{diag}(2,0)$，可是$f=x^2+y^2$在直角坐标中的Hessian是$2\bm I$，在所有方向上的二阶变化都相同。差别来自坐标方向场本身的变化：沿$\partial_\theta$对$\partial_\theta f$求导时，$\partial_\theta$也在转动。

\begin{definition}[Hessian]
\label{def:mcalc-hessian}
设$(M,g)$是黎曼流形，$\nabla$是其Levi-Civita联络，$f\in C^\infty(M)$。对$X,Y\in\mathfrak X(M)$定义
\begin{equation}
\operatorname{Hess}f(X,Y)=X(Yf)-df(\nabla_XY).
\label{eq:mcalc-hessian}
\end{equation}
\end{definition}

第一项是先沿$Y$求导再沿$X$求导，第二项扣除方向场$Y$沿$X$变化所带来的那部分。扣除后，$\operatorname{Hess}f$对$Y$按函数线性：$X(hYf)-df(\nabla_X(hY))=(Xh)Yf+hX(Yf)-(Xh)Yf-h\,df(\nabla_XY)$，交叉项抵消。对$X$的函数线性由联络的第一条规则直接得到。所以$\operatorname{Hess}f$是每点处$T_pM$上的双线性型，输入是函数，输出是双线性型，用到了联络。在坐标中，
\[
\operatorname{Hess}f(\partial_i,\partial_j)=\partial_i\partial_jf-\sum_k\Gamma^k_{i,j}\,\partial_kf .
\]
回到极坐标中的$f=r^2$：$\operatorname{Hess}f(\partial_r,\partial_r)=2$，$\operatorname{Hess}f(\partial_\theta,\partial_\theta)=0-\Gamma^r_{\theta,\theta}\cdot2r=2r^2$，交叉项为零，即$\operatorname{Hess}f=2g$，与直角坐标中的$2\bm I$一致。

\begin{theorem}[Hessian的性质]
\label{thm:mcalc-hessian-properties}
在Levi-Civita联络下：
\begin{enumerate}
\item $\operatorname{Hess}f(X,Y)=g\bigl(\nabla_X\operatorname{grad}f,\,Y\bigr)$；
\item $\operatorname{Hess}f$是对称双线性型：$\operatorname{Hess}f(X,Y)=\operatorname{Hess}f(Y,X)$；
\item 对任一测地线$\gamma$，$\dfrac{\mathrm d^2}{\mathrm dt^2}f(\gamma(t))=\operatorname{Hess}f\bigl(\dot\gamma(t),\dot\gamma(t)\bigr)$。
\end{enumerate}
\end{theorem}

\begin{proof}
(1) 由梯度的定义，$Yf=g(\operatorname{grad}f,Y)$。对它沿$X$求导并用度量相容，$X(Yf)=g(\nabla_X\operatorname{grad}f,Y)+g(\operatorname{grad}f,\nabla_XY)$，而最后一项是$df(\nabla_XY)$，移项即得。(2) 两种次序之差为$X(Yf)-Y(Xf)-df(\nabla_XY-\nabla_YX)=[X,Y]f-df([X,Y])=0$，其中用了无挠。(3) $\frac{\mathrm d}{\mathrm dt}f(\gamma)=df(\dot\gamma)=g(\operatorname{grad}f,\dot\gamma)$；再求导，由度量相容得$g(\nabla_{\dot\gamma}\operatorname{grad}f,\dot\gamma)+g(\operatorname{grad}f,D\dot\gamma/\mathrm dt)$，测地线使第二项为零，由(1)第一项就是$\operatorname{Hess}f(\dot\gamma,\dot\gamma)$。
\end{proof}

性质(3)给出Hessian的几何意义：它是函数沿测地线的二阶导数，正如欧氏Hessian是沿直线的二阶导数。性质(2)依赖无挠，对一般联络，式\eqref{eq:mcalc-hessian}仍有定义，却未必对称。还要注意，Hessian与第~\ref{sec:mcalc-forms}节的$2$形式都是两个切向量的双线性读数，但前者对称，后者交替，两类对象互不包含（除零之外）。把Hessian称为“二阶微分形式”会混淆这两者。

\subsection{散度}
\label{sec:mcalc-divergence}

把向量场$X$想成流体的速度，散度回答：在一点附近，流动使体积以多快的速率膨胀。欧氏空间中它是$\sum_i\partial_iX^i$，即Jacobian矩阵$[\partial_jX^i]$的迹。在流形上，Jacobian矩阵的角色由线性映射$v\mapsto\nabla_vX$承担。

\begin{definition}[散度]
\label{def:mcalc-divergence}
设$(M,g)$是黎曼流形，$X\in\mathfrak X(M)$。$X$的\term{散度}（divergence）是函数
\[
\operatorname{div}X(p)=\operatorname{tr}\bigl(T_pM\ni v\mapsto\nabla_vX\in T_pM\bigr).
\]
\end{definition}

在坐标中，$\nabla_{\partial_i}X=\sum_k(\partial_iX^k+\sum_j\Gamma^k_{i,j}X^j)\partial_k$，迹为$\sum_i\partial_iX^i+\sum_{i,j}\Gamma^i_{i,j}X^j$。对Levi-Civita联络，缩并的Christoffel符号有简单的形式：在第~\ref{chap:riemannian-geometry}章的公式$\Gamma^k_{i,j}=\frac12\sum_\ell g^{k,\ell}(\partial_ig_{j,\ell}+\partial_jg_{i,\ell}-\partial_\ell g_{i,j})$中令$k=i$并求和，第一项与第三项因$g^{i,\ell}$对称而抵消，剩下$\Gamma^i_{i,j}$求和为$\frac12\sum_{i,\ell}g^{i,\ell}\partial_jg_{i,\ell}=\frac12\operatorname{tr}(\bm G^{-1}\partial_j\bm G)$。由行列式的求导公式$\partial_j\det\bm G=\det\bm G\operatorname{tr}(\bm G^{-1}\partial_j\bm G)$，它等于$\partial_j\ln\sqrt{\det\bm G}$。代入后得到
\begin{equation}
\operatorname{div}X=\frac1{\sqrt{\det\bm G}}\sum_i\partial_i\bigl(\sqrt{\det\bm G}\,X^i\bigr).
\label{eq:mcalc-divergence-coordinates}
\end{equation}
输入是向量场，输出是函数。定义用的是联络，而坐标式只出现体积密度$\sqrt{\det\bm G}$：散度归根到底衡量的是黎曼体积的变化。

\paragraph{通量关系}
用内乘可以把散度与形式联系起来。对$k$形式$\eta$与向量场$X$，\term{内乘}（interior product）$\iota_X\eta$是把$X$填入第一个位置得到的$(k-1)$形式：$(\iota_X\eta)(v_1,\ldots,v_{k-1})=\eta(X,v_1,\ldots,v_{k-1})$。在定向流形上，对体积形式有
\begin{equation}
d(\iota_X\omega_g)=(\operatorname{div}X)\,\omega_g .
\label{eq:mcalc-divergence-form}
\end{equation}
验证只需在正向坐标图中计算：
\begin{equation*}
\begin{aligned}
&\iota_X(\sqrt{\det\bm G}\,dx^1\wedge\cdots\wedge dx^n)\\
&\qquad=\sum_i(-1)^{i-1}\sqrt{\det\bm G}\,X^i\,dx^1\wedge\cdots\wedge\widehat{dx^i}\wedge\cdots\wedge dx^n,
\end{aligned}
\end{equation*}
取外微分时第$i$项只剩$\partial_i(\sqrt{\det\bm G}X^i)$乘$dx^i$，移回第$i$位抵消符号，得到$\sum_i\partial_i(\sqrt{\det\bm G}X^i)\,dx^1\wedge\cdots\wedge dx^n$，与式\eqref{eq:mcalc-divergence-coordinates}相符。

\begin{theorem}[散度定理]
\label{thm:mcalc-divergence-theorem}
设$(M,g)$是定向的带边界黎曼流形，$X$是紧支撑的光滑向量场，$N$是边界上的外单位法向，$\mathrm dA$是边界上诱导度量的体积密度。则
\[
\int_M\operatorname{div}X\,\mathrm dV_g=\int_{\partial M}g(X,N)\,\mathrm dA .
\]
\end{theorem}

\begin{proof}
由式\eqref{eq:mcalc-divergence-form}与Stokes公式，左端等于$\int_{\partial M}\iota^*(\iota_X\omega_g)$。在边界点把$X$分解为$X=g(X,N)N+T$，$T$切于边界。$\iota^*(\iota_T\omega_g)$把$T$与$n-1$个边界切向量一起填入$\omega_g$，这$n$个向量都在$n-1$维的$T_p\partial M$中，线性相关，值为零。对$\iota^*(\iota_N\omega_g)$，取边界的正向标准正交基$(e_2,\ldots,e_n)$，由诱导定向$(N,e_2,\ldots,e_n)$是$T_pM$的正向标准正交基，故$\omega_g(N,e_2,\ldots,e_n)=1$；这说明$\iota^*(\iota_N\omega_g)$正是边界的体积形式。所以右端是$\int_{\partial M}g(X,N)\,\mathrm dA$。
\end{proof}

这就是上一节Gauss公式的流形形式，“面积元”现在有了准确的含义。定理不要求$M$可定向：两端都不依赖定向的选取，在不可定向时可以对局部可定向的坐标块分别应用再用单位分解拼接，这里不展开。另需注意，散度是关于黎曼体积定义的。若要关于另一个测度$q\,\mathrm dV_g$讨论通量平衡，相应的散度是$\frac1q\operatorname{div}(qX)$，二者相差$g(\operatorname{grad}\ln q,X)$，不能默认相同。

\subsection{Laplace--Beltrami算子}
\label{sec:mcalc-laplace-beltrami}

欧氏空间中，Laplace算子$\sum_i\partial_i^2f$是梯度的散度。流形上用同样的组合定义。

\begin{definition}[Laplace--Beltrami算子]
\label{def:mcalc-laplace-beltrami}
黎曼流形$(M,g)$上，对$f\in C^\infty(M)$定义
\begin{equation}
\Delta_gf=\operatorname{div}(\operatorname{grad}f)=\frac1{\sqrt{\det\bm G}}\sum_{i,j}\partial_i\Bigl(\sqrt{\det\bm G}\,g^{i,j}\,\partial_jf\Bigr).
\label{eq:mcalc-laplace-beltrami}
\end{equation}
$\Delta_g$称为\term{Laplace--Beltrami算子}（Laplace--Beltrami operator）。
\end{definition}

坐标式由式\eqref{eq:mcalc-gradient-coordinates}与\eqref{eq:mcalc-divergence-coordinates}合成。等价地，$\Delta_gf=\operatorname{tr}_g\operatorname{Hess}f$，即Hessian关于$g$的迹，这由定理~\ref{thm:mcalc-hessian-properties}(1)与散度的定义得到。输入是函数，输出是函数，度量通过$g^{i,j}$与$\sqrt{\det\bm G}$两处进入。符号约定须留意：本书的$\Delta_g$在欧氏空间中就是$\sum_i\partial_i^2$，下一节将看到它是“非正”的；不少文献取$-\operatorname{div}\operatorname{grad}$，使算子非负，读文献时要先确认约定。

\begin{example}[坐标中的Laplace--Beltrami算子]
\label{ex:mcalc-laplacian-examples}
欧氏直角坐标中$\bm G=\bm I$，式\eqref{eq:mcalc-laplace-beltrami}就是$\sum_i\partial_i^2f$。平面极坐标中$\sqrt{\det\bm G}=r$，$g^{r,r}=1$，$g^{\theta,\theta}=r^{-2}$，于是
\[
\Delta f=\frac1r\partial_r(r\,\partial_rf)+\frac1{r^2}\partial_\theta^2f=\partial_r^2f+\frac1r\partial_rf+\frac1{r^2}\partial_\theta^2f .
\]
单位球面上$\sqrt{\det\bm G}=\sin\theta$，
\[
\Delta_{\mathbb S^2}f=\frac1{\sin\theta}\partial_\theta(\sin\theta\,\partial_\theta f)+\frac1{\sin^2\theta}\partial_\phi^2f .
\]
取高度函数$f=\cos\theta$（即$z$），$\partial_\theta f=-\sin\theta$，$\Delta f=\frac1{\sin\theta}\partial_\theta(-\sin^2\theta)=-2\cos\theta$，所以$\Delta_{\mathbb S^2}z=-2z$。同一个函数$z$在$\mathbb R^3$中是线性函数，欧氏Laplacian为零；限制到球面上，Laplace--Beltrami算子却不为零。
\end{example}

\paragraph{依赖度量与离散近似}
Laplace--Beltrami算子随度量改变。二维情形有一个简单的规律：若$\tilde g=e^{2u}g$，则$\sqrt{\det\widetilde{\bm G}}=e^{2u}\sqrt{\det\bm G}$，$\tilde g^{i,j}=e^{-2u}g^{i,j}$，代入式\eqref{eq:mcalc-laplace-beltrami}，括号内的因子相消，得到$\Delta_{\tilde g}=e^{-2u}\Delta_g$；更高维时还会多出一阶项。在数据分析中，常用采样点上的图Laplacian近似$\Delta_g$；这种近似是否成立、收敛到哪个算子，取决于采样密度、核函数和尺度的选取，非均匀采样时极限往往含有密度带来的额外项。离散算子与几何算子是两个对象，二者的关系属于第~\ref{chap:manifold-recovery}章与卷九的内容。
